% ═══════════════════════════════════════════════════════════════════════════════
\section{Diskussion}
\label{sec:diskussion}
% ═══════════════════════════════════════════════════════════════════════════════

\subsection{Stärken des UQTL-Standards}
\label{subsec:staerken}

Der UQTL-Standard bietet gegenüber dem Status quo mehrere fundamentale
Vorteile:

\paragraph{Verifikationsgrundlage}
Durch das definierte Normalformat können Query-Äquivalenz\-prüfungen vollständig
automatisiert werden. Dies ermöglicht:
\begin{itemize}[leftmargin=1.5cm]
  \item Automatisierte Regressionstests bei Frameworkmigrationen
  \item Statische Analyse von Abfrage-Äquivalenz in CI/CD-Pipelines
  \item Formale Verifikation kritischer Datenzugriffspfade
\end{itemize}

\paragraph{Reduzierung kognitiver Last}
Entwickler, die in mehrsprachigen Teams arbeiten, können Abfragen in ihrer
bevorzugten Sprache formulieren und sind dennoch garantiert, dass semantisch
gleiche Abfragen zu identischen SQL-Anweisungen führen. Code-Reviews werden
deutlich vereinfacht, da der Vergleich auf der Ebene des normalisierten SQL
stattfinden kann.

\paragraph{Langzeitstabilität}
Frameworks evolvieren: SQLAlchemy hat mit Version 2.0 eine neue API eingeführt,
Entity Framework hat mehrere inkompatible Generationen durchlaufen. Ein
stabiler Standard abstrahiert von diesen Evolutionen.

\subsection{Grenzen und Limitierungen}
\label{subsec:grenzen}

\paragraph{Semantisch nicht äquivalente Quellkonstrukte}
Nicht alle Quellsprachkonstrukte haben eindeutige SQL-Entsprechungen.
Beispielsweise kann eine Python-Lambda-Funktion in einem \texttt{filter()}-Ausdruck
beliebige Python-Logik enthalten, die kein SQL-Äquivalent hat. UQTL definiert
für solche Fälle einen \emph{Fallback-Mechanismus}: Der betreffende Ausdruck
wird als \texttt{OPAQUE}-Knoten im IQT markiert und löst einen Normalisierungsfehler
aus.

\paragraph{Dialektspezifische Funktionen}
Bestimmte Datenbankfunktionen (z.\,B.\ PostgreSQL-spezifisches \texttt{jsonb\_agg()},
Oracle-spezifisches \texttt{CONNECT BY}) sind außerhalb von Erweiterungen nicht
im Standard enthalten. Implementierungen, die diese benötigen, müssen einen
Dialekt-Erweiterungsmechanismus implementieren.

\paragraph{Performanz-Implikationen}
Die Normalisierung erzeugt in einigen Fällen SQL-Ausdrücke, die semantisch
korrekt, aber nicht optimal für alle Datenbankplaner sind. Implementierungen
sind ermutigt, nach der Normalisierung einen optionalen \emph{Hint-Layer}
anzuwenden.

\paragraph{Stream-API vs. Datenbankabfragen}
Java Streams operieren in-memory und unterstützen Operationen (z.\,B.\
\texttt{peek()}, \texttt{forEach()}, zustandsbehaftete Lambdas), die keine
SQL-Entsprechung haben. UQTL richtet sich an den Subset der Stream-Operationen,
der eine relationale Semantik trägt.

\subsection{Erweiterungen und zukünftige Arbeit}
\label{subsec:erweiterungen}

Folgende Erweiterungen sind für zukünftige Versionen des Standards geplant:

\begin{enumerate}[leftmargin=1.5cm]
  \item \textbf{UQTL-E1 (NoSQL-Dialekte):} Erweiterung um MongoDB-Aggregations\-pipeline
    und Elasticsearch-DSL als zusätzliche Quellsprachen
  \item \textbf{UQTL-E2 (TypeScript):} Unterstützung von TypeScript/Prisma als
    vierte Quellsprache
  \item \textbf{UQTL-E3 (Streaming SQL):} Integration von Apache Flink SQL und
    Apache Spark SQL als Zieldialekte
  \item \textbf{UQTL-E4 (Schema-Awareness):} Integration von Schemavalidierung
    und Typ-Prüfung in den Normalisierungsprozess
  \item \textbf{UQTL-E5 (ML-gestützte Optimierung):} Automatische Query-Optimierung
    basierend auf historischen Ausführungsstatistiken
\end{enumerate}

% ═══════════════════════════════════════════════════════════════════════════════
\section{Zusammenfassung}
\label{sec:zusammenfassung}
% ═══════════════════════════════════════════════════════════════════════════════

Die vorliegende Arbeit hat den \textbf{Unified Query Translation Layer (UQTL)}
vollständig spezifiziert: einen offenen Standard zur garantierten semantischen
Äquivalenz von Query-Übersetzungen aus Python, Java (Stream API / JPA) und
C\#/.NET (LINQ) nach SQL.

\paragraph{Kernbeiträge}
Die wesentlichen Beiträge des Standards lassen sich wie folgt zusammenfassen:

\begin{enumerate}[leftmargin=1.5cm]
  \item \textbf{Formale Grundlage:} Die Definition des Intermediate Query Tree
    (IQT) als typisierter DAG mit vollständigen Knotenspezifikationen für alle
    11~Knotentypen (SCAN, FILTER, PROJECT, JOIN, GROUP\_AGG, SORT, LIMIT,
    WINDOW, SUBQUERY, SET\_OP, WITH).

  \item \textbf{Vollständige Übersetzungsregeln:} Für alle 9~Hauptoperatoren
    (Filter, Projektion, Sortierung, Limit, Aggregation, Joins, Fensterfunktionen,
    Unterabfragen, Mengenoperationen) wurden vollständige, formale
    Ãœbersetzungsregeln von allen drei Quellsprachen in den IQT und vom IQT
    nach ANSI SQL:2016 definiert.

  \item \textbf{10~Praxisbeispiele:} Anhand eines einheitlichen E-Commerce-Daten\-modells
    wurden 10~vollständige Beispiele demonstriert, die von einfacher Filterung bis
    zu komplexen Geschäftsberichten mit CTEs und Fensterfunktionen reichen.

  \item \textbf{Konformitätsstufen:} Vier Stufen (CL-1 bis CL-4) erlauben
    einen schrittweisen Aufbau konformer Implementierungen.

  \item \textbf{Referenzimplementierungsrichtlinien:} Architektur und
    Pseudocode einer konformen Implementierung wurden bereitgestellt.
\end{enumerate}

\paragraph{Bedeutung für die Praxis}

Der UQTL-Standard schließt eine fundamentale Lücke in der polyglotten
Softwareentwicklung: Erstmals können Teams aus Python-, Java- und
C\#-Entwicklern formal garantieren, dass ihre Datenbankzugriffe semantisch
äquivalent sind -- unabhängig von Sprache, Framework-Version oder persönlichem
Codingstil.

\begin{successbox}
\textbf{Leitprinzip des Standards:}

\emph{Dieselbe Datenabfrage, formuliert in unterschiedlichen Sprachen,
erzeugt unter UQTL stets dieselbe SQL-Anweisung. Die Sprache ist ein
Implementierungsdetail -- die Semantik ist der Standard.}
\end{successbox}

